\section{OrdMon (Monoides Ordenados)}

A categoria \textbf{OrdMon} une álgebra e ordem. Um monoide ordenado exige que a operação binária seja compatível com a relação de ordem, ou seja, se $a \le b$, então $a*c \le b*c$.

\subsection{Objetos}
A representação destes objetos requer uma estrutura híbrida que combine a Tabela de Cayley (para a parte algébrica) com uma matriz de adjacência (para a componente de ordem).

\begin{lstlisting}[caption={Modelação Híbrida de OrdMon}]
OM.elements = {1, 2, 3};
OM.mul = @(a, b) a + b; 
OM.unit = 1;            
OM.le = @(a, b) a <= b; 
\end{lstlisting}

\subsection{Morfismos}
Os morfismos em \textbf{OrdMon} são os que satisfazem simultaneamente as condições de homomorfismo de monoides e de monotonia de ordens parciais.